%----------------------------------------------------------------------------------------
% HAAKJES UITWERKEN EN MERKWAARDIGE PRODUCTEN. Alle opgaven zijn eigen werk.
%
% Een hoofdstuk kernvaardigheden (ontwerpregels, deel Kernvaardigheden): beknopte theorie, daarna drill.
% Labels van dit hoofdstuk beginnen met h2-.
%----------------------------------------------------------------------------------------

\mychapter{2}{Merkwaardige producten}

\lkr{Herhaling van de 1ste graad (rekenen met letters, merkwaardige producten). Nodig in Tweedegraadsfuncties (de topvorm uitwerken, kwadraat afsplitsen in bundel 8), Veeltermen en Analytische meetkunde (de vergelijking van een cirkel). Laat leerlingen de fout $(a + b)^2 = a^2 + b^2$ met getallen weerleggen: $a = b = 1$.}

\section{Wat je moet kennen}

\begin{theorie}
\begin{minipage}[t]{.47\linewidth}
\vspace{0pt}
\textbf{\textcolor{accent}{Distributiviteit}}

\vspace{1mm}
$\formule{a(b + c) = ab + ac}$

\vspace{1mm}
$\formule{(a + b)(c + d) = ac + ad + bc + bd}$
\end{minipage}\hfill
\begin{minipage}[t]{.47\linewidth}
\vspace{0pt}
\textbf{\textcolor{accent}{Merkwaardige producten}}

\vspace{1mm}
$\formule{(A + B)^2 = A^2 + 2AB + B^2}$

\vspace{1mm}
$\formule{(A - B)^2 = A^2 - 2AB + B^2}$

\vspace{1mm}
$\formule{(A + B)(A - B) = A^2 - B^2}$
\end{minipage}
\end{theorie}

\textbf{Klassieke fouten.} Het dubbel product vergeten: $(x + 3)^2 \neq x^2 + 9$. Een factor niet mee kwadrateren: $(3a)^2 = 9a^2$, niet $3a^2$. Een minteken voor haakjes geldt voor elke term: $-(x - 2)^2 = -(x^2 - 4x + 4) = -x^2 + 4x - 4$. Werk daarom eerst elk product uit \emph{tussen haakjes}, en haal pas daarna de haakjes weg.

\voorbeeld{}{Vereenvoudig $(x + 3)^2 - (x - 2)(x + 2) - 2(x - 1)^2$.\label{vb:h2-gemengd}}{%
\begin{stappen}
(x^2 + 6x + 9) - (x^2 - 4) - 2(x^2 - 2x + 1) & Elk product eerst tussen haakjes. \\
= x^2 + 6x + 9 - x^2 + 4 - 2x^2 + 4x - 2 & Haakjes weg: elk teken wisselt na een minteken. \\
= -2x^2 + 10x + 11 & Samennemen.
\end{stappen}}

\voorbeeld{}{Bereken uit het hoofd: \quad a. $98^2$ \qquad b. $103 \cdot 97$\label{vb:h2-handig}}{%
\begin{stappen}
(100 - 2)^2 = 10\,000 - 400 + 4 = 9604 & a. Kwadraat van een verschil. \\
(100 + 3)(100 - 3) = 10\,000 - 9 = 9991 & b. Toegevoegde tweetermen.
\end{stappen}}

%========================================================================================
\opdrachtenkop

\begin{oefeningen}
\needspace{10\baselineskip}
\oef{h2-o1}

\needspace{10\baselineskip}
\oef{h2-o2}

\needspace{10\baselineskip}
\oef{h2-o3}

\needspace{10\baselineskip}
\oef{h2-o4}

\needspace{10\baselineskip}
\oef{h2-o5}

\needspace{10\baselineskip}
\oef{h2-o6}

\needspace{10\baselineskip}
\oef{h2-o7}

\needspace{10\baselineskip}
\oef{h2-o8}

\needspace{10\baselineskip}
\oef{h2-o9}

\needspace{10\baselineskip}
\oef{h2-o10}
\end{oefeningen}

%----------------------------------------------------------------------------------------
\needspace{14\baselineskip}
\opdrachtsectie{Test jezelf}

\begin{oefeningen}
\oef{h2-t1}

\oef{h2-t2}

\oef{h2-t3}

\oef{h2-t4}

\oef{h2-t5}
\end{oefeningen}

\herhaalschema{%
  \fv{oef:h2-t1}{vb:h2-gemengd}{oef:h2-o1}%
  \fzv{oef:h2-t2}{oef:h2-o2}%
  \fzv{oef:h2-t3}{oef:h2-o3}%
  \fv{oef:h2-t4}{vb:h2-gemengd}{oef:h2-o6}%
  \fv{oef:h2-t5}{vb:h2-handig}{oef:h2-o4}}

\ifzelfstudie
\section*{Maak je eigen samenvatting}
Sluit dit boek en schrijf in je schrift: de distributiviteit, de drie merkwaardige producten, en waarom $(A + B)^2$ niet gelijk is aan $A^2 + B^2$ (met een voorbeeld met getallen). Vergelijk daarna met het kader.
\fi
